\subsection*{Formulação do Modelo}
O problema descreve a queda de um corpo de massa $m$ sob a influência da gravidade $g$ e de uma força de resistência do ar proporcional ao quadrado da velocidade ($F_r = cv^2$). Diferente do modelo básico, este sistema possui uma \textbf{velocidade inicial} $v(0) = v_0$.

A Segunda Lei de Newton estabelece a seguinte Equação Diferencial Ordinária (EDO):
\begin{equation}
    m \frac{dv}{dt} = mg - cv^2
\end{equation}

Definimos a velocidade terminal ($v_t$) como o estado onde a aceleração é nula ($mg = cv_t^2$). Assim, podemos expressar a constante de arrasto como $c = \frac{mg}{v_t^2}$. Substituindo na EDO original, obtemos a forma normalizada:
\begin{equation}
    \frac{dv}{dt} = g \left( 1 - \frac{v^2}{v_t^2} \right)
\end{equation}

\subsection*{Separação de Variáveis e Integração}
Para determinar o tempo $t$ necessário para atingir uma velocidade $v$, isolamos os termos e integramos nos limites definidos pela condição inicial:
\begin{equation}
    \int_{v_0}^{v} \frac{dv}{1 - \left( \frac{v}{v_t} \right)^2} = \int_{0}^{t} g \, dt
\end{equation}

Utilizando a expansão por frações parciais para o termo à esquerda:
\begin{equation}
    \frac{1}{1 - (v/v_t)^2} = \frac{v_t^2}{v_t^2 - v^2} = \frac{v_t}{2} \left( \frac{1}{v_t + v} + \frac{1}{v_t - v} \right)
\end{equation}

Integrando ambos os lados:
\begin{equation}
    \frac{v_t}{2} \left[ \ln \left( \frac{v_t + v}{v_t - v} \right) \right]_{v_0}^{v} = gt
\end{equation}

\subsection*{Solução Analítica}
Aplicando os limites de integração superiores e inferiores, obtemos a expressão para o tempo $t$:
\begin{equation}
    t = \frac{v_t}{2g} \left[ \ln \left( \frac{v_t + v}{v_t - v} \right) - \ln \left( \frac{v_t + v_0}{v_t - v_0} \right) \right]
\end{equation}

Utilizando as propriedades de logaritmo ($\ln A - \ln B = \ln \frac{A}{B}$), a solução final compacta é:
\begin{equation}
    t = \frac{v_t}{2g} \ln \left( \frac{(v_t + v)(v_t - v_0)}{(v_t - v)(v_t + v_0)} \right)
\end{equation}

\subsection*{Análise de Convergência}
Seja $\alpha$ a fração da velocidade terminal que se deseja atingir ($v = \alpha v_t$, com $0 < \alpha < 1$). O tempo $t_\alpha$ é dado por:
\begin{equation}
    t_{\alpha} = \frac{v_t}{2g} \ln \left( \frac{(1 + \alpha)(v_t - v_0)}{(1 - \alpha)(v_t + v_0)} \right)
\end{equation}

\textbf{Conclusão:} A presença da velocidade inicial $v_0$ atua como uma translação na escala temporal. Se $v_0 = 0$, a solução converge para o modelo de queda livre padrão. Se $v_0 > v_t$, o objeto desacelerará até atingir a velocidade terminal, mantendo a validade física da solução logarítmica (em módulo).
